% Schema COMPLET du modele : donnees, calibration, mecanisme, identification, sorties.
% Compile avec Tectonic. Une page A4 paysage, incluable dans le memoire et les slides.
%
% PRINCIPE DE LECTURE, ET C'EST LE POINT DU SCHEMA. Chaque grandeur porte une pastille qui
% dit son STATUT epistemique : estimee sur donnee, posee sans calibration, bornee faute
% d'identification, ou reglementaire. Un lecteur doit pouvoir attaquer le modele bloc par
% bloc, et savoir en un coup d'oeil ce qui se discute et ce qui se verifie.
\documentclass[11pt,a4paper,landscape]{article}
\usepackage[utf8]{inputenc}
\usepackage[T1]{fontenc}
\usepackage{lmodern}
\usepackage[french]{babel}
\usepackage[margin=0.38cm]{geometry}
\usepackage{amsmath,amssymb}
\usepackage{eurosym}
\usepackage[table]{xcolor}
\usepackage{tikz}
\usetikzlibrary{positioning,arrows.meta,fit,backgrounds,calc,shapes.geometric}

\definecolor{navy}{HTML}{1F3864}
\definecolor{bluemid}{HTML}{2E5496}
\definecolor{lightblue}{HTML}{D6E0F0}
\definecolor{accent}{HTML}{C0491F}
\definecolor{vert}{HTML}{3D8361}
\definecolor{ocre}{HTML}{B07D2B}
\definecolor{gris}{HTML}{6B6B6B}
\definecolor{grisclair}{HTML}{ECECEC}

\pagestyle{empty}

\tikzset{
  bloc/.style={
    draw=navy, fill=lightblue!25, line width=0.6pt, rounded corners=2pt,
    align=left, inner sep=3.5pt, font=\scriptsize, text width=#1},
  bloc/.default=8.7cm,
  sortie/.style={
    draw=vert!80!black, fill=vert!8, line width=0.7pt, rounded corners=2pt,
    align=left, inner sep=3.5pt, font=\scriptsize, text width=#1},
  sortie/.default=8.7cm,
  limite/.style={
    draw=accent!70!black, fill=accent!6, line width=0.7pt, rounded corners=2pt,
    align=left, inner sep=3.5pt, font=\scriptsize, text width=#1},
  limite/.default=8.7cm,
  entete/.style={
    fill=navy, text=white, rounded corners=2pt, align=center,
    inner sep=2.5pt, font=\small\bfseries, text width=8.7cm},
  pilier/.style={
    draw=bluemid, fill=white, line width=0.7pt, rounded corners=2pt,
    align=center, inner sep=1.5pt, font=\tiny, text width=1.34cm, minimum height=0.62cm},
  fl/.style={-{Stealth[length=2.0mm]}, line width=0.7pt, color=navy},
  flfin/.style={-{Stealth[length=1.5mm]}, line width=0.5pt, color=bluemid},
  flcol/.style={-{Stealth[length=2.4mm]}, line width=1.0pt, color=navy!70},
}

% --- pastilles de statut : le coeur de la lecture
\newcommand{\pest}{\textcolor{bluemid}{$\blacksquare$}}   % estime sur donnee
\newcommand{\ppos}{\textcolor{ocre}{$\blacklozenge$}}     % pose, non calibre
\newcommand{\pbor}{\textcolor{accent}{$\blacktriangle$}}  % borne, non identifie
\newcommand{\preg}{\textcolor{gris}{$\bullet$}}           % reglementaire ou standard

\newcommand{\eq}[1]{\par\vspace{1.5pt}{\centering$\displaystyle #1$\par}\vspace{1.5pt}}
\newcommand{\ttl}[1]{{\color{navy}\bfseries #1}\\[0.5pt]}
\newcommand{\src}[1]{\hfill{\tiny\color{gris}script~#1}}

\begin{document}
\begin{center}

{\color{navy}\large\bfseries Quantifier le besoin de capital DORA : de la donnée brute au
capital, puis à son allocation}\\[1pt]
{\scriptsize\color{gris}
\pest~estimé sur donnée \quad
\ppos~posé, non calibré \quad
\pbor~borné faute d'identification \quad
\preg~réglementaire ou standard
\;\textbar\; chaque bloc est un choix attaquable séparément, et chaque nombre est imprimé par
le script indiqué à droite.}

\vspace{0.8mm}

\begin{tikzpicture}[node distance=1.7mm]

% =====================================================================
% COLONNE 1 : DONNEES ET CALIBRATION
% =====================================================================
\node[entete] (h1) {I. Les données, et ce qu'on en tire};

\node[bloc, below=of h1] (b11)
  {\ttl{1. Trois sources, trois rôles disjoints \src{62, 63}}
   \pest~\textbf{OpRisk Global} (SAS), $582$ pertes cyber${\times}$finance sur $27$ ans,
   converties en euros $\to$ porte le \emph{niveau} de sévérité.\\
   \pest~\textbf{Privacy Rights Clearinghouse}, $15\,053$ incidents en nombre
   d'enregistrements, convertis par la régression de Jacobs
   $\ln L = 7{,}68+0{,}76\ln X+\varepsilon$ $\to$ axe de \emph{sensibilité}.\\
   \pest~\textbf{Corpus de post-mortems}, $10$ incidents et $22$ transitions codées
   $\to$ la seule source d'\emph{ordre} entre piliers.\\[1pt]
   \preg~\textbf{Hackmageddon} : parts par vecteur seulement, citation externe datée et
   \emph{non recalculable}, qui ne porte aucun niveau de capital.};

\node[bloc, below=of b11] (b12)
  {\ttl{2. Fréquence : une charge systémique commune \src{08b, 41}}
   $Y\sim\mathcal{N}(0,1)$ tirée \emph{une fois par an}, commune aux cinq piliers~:
   \eq{
     m=\exp\!\bigl(aY-\tfrac{a^2}{2}\bigr),\quad \mathbb{E}[m]=1,\qquad
     N_j \mid Y \sim \mathcal{P}\bigl(\lambda_j\, m\bigr)
   }
   \ppos~$a=0{,}60$ \emph{posé}, contre \pest~$a=0{,}122$ \emph{estimé}
   (IC$_{95\%}$ $[0{,}054\,;0{,}223]$) : \textbf{le mémoire retient donc une charge
   systémique plus lourde que sa propre estimation}, et le dit.\\
   \pest~$\lambda_{\text{secteur}}=21{,}56$/an, \; $\lambda_j\propto\mathrm{ROOT}_j$,
   parts $30/27/18/15/9\,\%$. Surdispersion observée $\mathrm{Var}/\mathbb{E}=1{,}19$,
   dispersion de Pearson $\varphi=2{,}24$.};

\node[bloc, below=of b12] (b13)
  {\ttl{3. Sévérité : loi recollée, corps empirique et queue GPD \src{07, 47, 63}}
   \eq{
     \mathbb{P}(X>x)=p_u\Bigl(1+\tfrac{\xi(x-u)}{\sigma}\Bigr)^{-1/\xi},\qquad x>u
   }
   \pest~\textbf{OpRisk} : $u=20{,}03$~M\euro, $91$ excès, $\xi=0{,}595$
   (IC$_{90\%}$ $[0{,}30\,;0{,}83]$), $\sigma=57{,}97$~M\euro, $p_u=0{,}151$.\\
   \pest~\textbf{PRC} : $u=4{,}176$~M\euro, $2\,258$ excès, $\xi=1{,}033$ :
   \emph{espérance infinie}, ce qui justifie \ppos~un plafond de sévérité à $40$~M\euro.\\[1pt]
   \textbf{Validé, pas supposé} : Anderson-Darling $p=0{,}99$ et Kolmogorov-Smirnov
   $p=0{,}89$ sur les paramètres \emph{publiés}, paramètres imposés. $\xi$ ne bouge que de
   $0{,}5\,\%$ si l'on déplace le seuil au percentile $85$.};

\node[limite, below=of b13] (b14)
  {\ttl{4. Ce que la donnée ne permet pas, et qu'on ne masque pas \src{57, 67}}
   \pbor~Biais de taille : OpRisk sur-représente les grandes institutions, la queue est
   possiblement sous-estimée. Biais de troncature au seuil de collecte.\\
   \pbor~Aucune base ne date les enchaînements : $4\,329$ transitions annuelles ne donnent
   \emph{aucun} signal directionnel ($z=-0{,}33$).\\
   \pbor~Le second codage en aveugle du corpus reste à faire : la fiabilité inter-juges
   n'est pas mesurée.};

% =====================================================================
% COLONNE 2 : MECANISME ET IDENTIFICATION
% =====================================================================
\node[entete, right=6mm of h1] (h2) {II. Le mécanisme, et sa part inidentifiable};

\node[bloc, below=of h2] (b21)
  {\ttl{5. De la matrice de jugement à la matrice de contagion \src{16, 20, 30}}
   $s_j=\sum_k \mathrm{TRANS}_{jk}$ est la \emph{puissance de propagation} du pilier $j$~:
   \eq{
     W_{jk}=g\,\frac{\mathrm{TRANS}_{jk}}{\max_\ell s_\ell},\qquad
     e_j = g\,\frac{s_j}{\max_\ell s_\ell}
   }
   \ppos~$g$ par état de conformité : $0{,}45$ / $0{,}68$ / $0{,}90$, \emph{posés}.\\
   \pest~$\rho(W)=0{,}506$ à $g=0{,}90$. Comme $\sum_k W_{jk}\le g$, on a $\rho(W)\le g<1$~:
   \textbf{sous-criticité garantie par construction}, la cascade s'éteint toujours.};

\node[pilier, below=5mm of b21.south west, xshift=0.88cm] (p1) {\textbf{P1}\\gouvernance};
\node[pilier, right=2.2mm of p1] (p2) {\textbf{P2}\\incidents};
\node[pilier, right=2.2mm of p2] (p3) {\textbf{P3}\\tests};
\node[pilier, right=2.2mm of p3] (p4) {\textbf{P4}\\tiers};
\node[pilier, right=2.2mm of p4] (p5) {\textbf{P5}\\partage};

\draw[flfin] (p1.south) to[bend right=22] (p2.south);
\draw[flfin] (p4.south) to[bend left=18] (p2.south);
\draw[flfin] (p2.south) to[bend right=16] (p5.south);
\draw[flfin] (p5.south) to[bend left=20] (p3.south);
\draw[flfin] (p1.south) to[bend right=13] (p3.south);
\node[font=\tiny\itshape, color=gris] (lbr) at ($(p3.south)+(0,-0.82)$)
  {les cinq courants nets dominants ; $\max_{jk}|J_{jk}|=0{,}044$, aucun n'est significatif};

\begin{scope}[on background layer]
  \node[draw=navy, dashed, line width=0.5pt, rounded corners=3pt,
        fit=(p1)(p2)(p3)(p4)(p5)(lbr), inner sep=3pt, fill=white] (casc) {};
\end{scope}

\node[bloc, below=2mm of casc] (b22)
  {\ttl{6. Le branchement, et la loi exacte de l'étendue \src{16b}}
   Depuis un pilier $j$ tombé, chaque cible est tirée proportionnellement à
   $\mathrm{TRANS}_{jk}/s_j$ et infectée avec probabilité $W_{jk}$. La loi du
   \textbf{nombre de piliers touchés} est calculée \emph{exactement}, sans Monte-Carlo,
   par récursion sur les $2^5$ sous-ensembles~:
   \eq{
     \mathbb{P}(\mathcal{S}=\mathcal{A})=R(\mathcal{A})
     \!\!\prod_{j\in\mathcal{A},\,k\notin\mathcal{A}}\!\!(1-W_{jk})
   }
   où le produit est la probabilité qu'\emph{aucune} fuite ne sorte de $\mathcal{A}$.};

\node[limite, below=of b22] (b23)
  {\ttl{7. Le c\oe{}ur du problème : la direction n'est pas identifiée \src{40, 53, 59}}
   Décomposition $W=S+A$, avec $S$ symétrique (\emph{le lien}) et $A$ antisymétrique
   (\emph{le sens}).\\
   \pest~$S$ est identifié : la co-occurrence se lit dans les données.\\
   \pbor~$A$ ne l'est pas : trois tests successifs le montrent.
   Placebo par permutation sur le pas annuel d'OpRisk, $z=-0{,}33$, \emph{aucun signal} ;
   sur les post-mortems, $z=+3{,}93$ contre un nul de $8{,}4\pm2{,}2$ ; sur le corpus
   étendu, $z=+5{,}12$.\\[1pt]
   \textbf{Réponse : l'identification partielle.} On énumère \emph{exhaustivement} les
   $2^{10}=1024$ sommets de la boîte admissible $|A_{jk}|\le t\,S_{jk}$, et l'on rapporte
   une \textbf{bande}, pas un point.};

\node[bloc, below=of b23] (b24)
  {\ttl{8. Amplification, et adaptations par pilier \src{37, 38, 39}}
   $(I-W)^{-1}=I+W+W^2+\cdots$ converge car $\rho(W)<1$. \emph{La direction n'apparaît
   qu'à partir de $W^2$ : c'est là que se joue la dépendance à l'ordre.}\\[1pt]
   \pbor~Queue par pilier : cinq queues Bâle observées, $0{,}87$ à $1{,}37$, non
   assignables aux piliers. Les $120$ assignations bornent le SCR à $+55$ / $+96\,\%$.\\
   \pest~P4 (tiers) n'est pas un n\oe{}ud comme les autres : $191$ victimes d'un seul choc
   (MOVEit) $\to$ \textbf{choc commun}, structure symétrique donc \emph{identifiée}.};

% =====================================================================
% COLONNE 3 : SORTIES, ECHELLE, LIMITES
% =====================================================================
\node[entete, right=6mm of h2] (h3) {III. Ce qui sort, et à quelles conditions};

\node[sortie, below=of h3] (b31)
  {\ttl{9. Agrégation et capital \src{14, 16, 21}}
   Modèle collectif, $L=\sum_{i=1}^{N} X_i^{\text{amplifiée}}=\sum_{j=1}^{5} L_j$,
   $60\,000$ à $600\,000$ années simulées, nombres aléatoires communs entre configurations.\\
   \preg~$\mathrm{SCR}=\mathrm{VaR}_{99{,}5\,\%}(L)$, forme imposée par Solvabilité~II.\\[1pt]
   socle sans contagion \; $5\,275$~M\euro \\
   \pbor~avec contagion, ignorance directionnelle totale \;
   $[\,6\,858\;;\;8\,697\,]$~M\euro\\
   surcoût de non-conformité \; $\Delta_{\text{DORA}}=10\,927$~M\euro};

\node[sortie, below=of b31] (b32)
  {\ttl{10. Allocation aux cinq piliers : deux vues, deux questions \src{20}}
   \textbf{Shapley}, vue \emph{source}, pour le surcoût~:
   $\varphi_j=\sum_{\mathcal{A}\not\ni j}
   \frac{|\mathcal{A}|!\,(4-|\mathcal{A}|)!}{5!}
   [v(\mathcal{A}\cup\{j\})-v(\mathcal{A})]$, avec
   $v(\mathcal{A})=\mathrm{SCR}(\mathcal{A}\text{ non conformes})-\mathrm{SCR}(\text{tous
   conformes})$. P1 concentre $34\,\%$, au-dessus de sa part d'amorce de $30\,\%$~:
   \emph{la propagation déforme l'amorce}.\\
   \textbf{Euler}, vue \emph{réceptacle}, pour le niveau~:
   $C_j=\mathbb{E}[L_j \mid L\ge \mathrm{VaR}_{99{,}5\,\%}]$, presque plate.\\[1pt]
   \emph{La remédiation se décide côté source, le capital se loge côté réceptacle.}};

\node[sortie, below=of b32] (b33)
  {\ttl{11. Descente à l'échelle d'une entité, et emploi \src{58, 60, 65}}
   $\lambda$ lu \emph{à la taille} de l'entité, $0{,}092$/an, et non choisi dans un seau.
   Besoin de capital $169$~M\euro, bornes $[\,131{,}5\;;\;183{,}3\,]$.\\
   \textbf{Arbitrage détenir ou transférer} : le coût annuel du capital vaut $2{,}2$ fois
   la sinistralité attendue ; la portée optimale suit le capital, et la conformité fait
   tomber le coût total de $2{,}89$ à $0{,}54$~M\euro/an, soit $-81\,\%$.\\
   \pest~Quatre bilans SFCR réels, entités anonymisées : le rapport au SCR publié est une
   \textbf{mise à l'échelle, pas une part}.};

\node[limite, below=of b33] (b34)
  {\ttl{12. Les réserves, publiées et non dissimulées \src{46, 48, 51, 66}}
   \pbor~\textbf{Trois étages d'incertitude} : paramètre $\times 2{,}5$ \; $\cdot$ \;
   dépendance $\times 1{,}8$ \; $\cdot$ \; famille de queue $\times 5{,}5$.
   \textbf{Le niveau bouge dans la bande, l'ordre des piliers ne bouge pas.}\\
   \pbor~Le $169$~M\euro{} est une \textbf{borne supérieure d'ordre de grandeur}, pas une
   mesure : l'entité notionnelle tombe sous la borne inférieure de validité de la descente.\\
   \ppos~Les trois $g$ sont posés, mais la thèse n'en dépend pas : le capital est
   \emph{croissant} en $g$ (lemme vérifié), donc tout ordre respecté produit l'écart.
   Seule l'\textbf{amplitude} est un scénario.\\
   \preg~\textbf{Ce n'est pas un SCR réglementaire} : aucun module DORA n'existe en Formule
   Standard. On écrit \og{}besoin de capital ORSA au titre de DORA\fg{}.};

% =====================================================================
% fleches
% =====================================================================
\draw[fl] (b11.south) -- (b12.north);
\draw[fl] (b12.south) -- (b13.north);
\draw[fl] (b13.south) -- (b14.north);
\draw[fl] (b21.south) -- (casc.north);
\draw[fl] (casc.south) -- (b22.north);
\draw[fl] (b22.south) -- (b23.north);
\draw[fl] (b23.south) -- (b24.north);
\draw[fl] (b31.south) -- (b32.north);
\draw[fl] (b32.south) -- (b33.north);
\draw[fl] (b33.south) -- (b34.north);

\draw[flcol] ($(b12.east)+(0.10,0)$) -- ($(b12.east)+(0.50,0)$);
\draw[flcol] ($(b21.east)+(0.10,0)$) -- ($(b21.east)+(0.50,0)$);
\draw[flcol] ($(b24.east)+(0.10,0)$) -- ($(b24.east)+(0.50,0)$);

\end{tikzpicture}

\vspace{0.4mm}
{\scriptsize\color{gris}
\textbf{Le fil, en une phrase} : la donnée identifie le \emph{lien} entre piliers, pas son
\emph{sens} ; on ne pose donc pas la direction, on l'encadre, et le classement des piliers
survit à cet encadrement là où le niveau du capital, lui, n'y survit pas.}

\end{center}
\end{document}
